\section{When Does Fixed Capacity Cause Forgetting?}
\label{sec:theory}

We use one adapted matrix when discussing rank and a vector $\bx$ for local
optimization geometry. Proofs are in Appendix~\ref{app:proofs}. All claims are
conditional on their stated smoothness, stationarity, curvature, and
block-structure assumptions.

\subsection{The Linear Term Disappears Only at the Right Stationary Point}

\begin{lemma}[Local expansion of signed forgetting]
\label{lem:local-forgetting}
Fix $s<t$ and let $\bm{d}_{s,t}:=\bx_t-\bx_s$. Suppose $\cL_s$ is twice
continuously differentiable and its Hessian is $\rho_s$-Lipschitz on the line
segment joining $\bx_s$ and $\bx_t$. With
$\bg_s=\nabla\cL_s(\bx_s)$ and $\bH_s=\nabla^2\cL_s(\bx_s)$,
\begin{equation}
  \cF_{s\to t}
  =\bg_s^\top\bm{d}_{s,t}
  +\frac12\bm{d}_{s,t}^\top\bH_s\bm{d}_{s,t}+r_{s,t},
  \qquad
  |r_{s,t}|\leq\frac{\rho_s}{6}\norm{\bm{d}_{s,t}}_2^3.
  \label{eq:forgetting-expansion}
\end{equation}
For a linear feasible update space $\cU$, the linear term vanishes for every
$\bm{d}_{s,t}\in\cU$ if and only if
$\Pi_{\cU}\bg_s=\bm{0}$.
\end{lemma}

Thus the boss-level observation that ``the linear term disappears and the
quadratic remains'' is correct only at projected stationarity for the same old
risk being evaluated. It does not hold automatically at an intermediate
checkpoint or for a shifted evaluation distribution. The quadratic term is a
nonnegative forgetting cost only when $\bH_s\succeq0$ on the relevant
directions.

\subsection{Sequential Updates Accumulate Only under Non-Cancellation}

Let $\bh_t:=\bx_t-\bx_{t-1}$.

\begin{theorem}[Conditional accumulation of sequential forgetting]
\label{thm:sequential-accumulation}
Fix a past window $s$ and suppose that along the later optimization path
\begin{equation}
  \cL_s(\bx)=\cL_s(\bx_s)
  +\frac12\norm{\bx-\bx_s}_{\bH_s}^2,
  \qquad \bH_s\succeq\bm{0}.
  \label{eq:quadratic-old-loss}
\end{equation}
Then
\begin{align}
  \cF_{s\to t}
  &=\frac12\norm{\sum_{u=s+1}^{t}\bh_u}_{\bH_s}^2,
  \label{eq:path-forgetting}\\
  \cF_{s\to t}-\cF_{s\to t-1}
  &=(\bx_{t-1}-\bx_s)^\top\bH_s\bh_t
  +\frac12\norm{\bh_t}_{\bH_s}^2.
  \label{eq:forgetting-increment}
\end{align}
If
\begin{equation}
  \E[(\bx_{t-1}-\bx_s)^\top\bH_s\bh_t]\geq-\beta_t,
  \label{eq:noncancellation}
\end{equation}
then
\begin{equation}
  \E[\cF_{s\to t}-\cF_{s\to t-1}]
  \geq\frac12\E\norm{\bh_t}_{\bH_s}^2-\beta_t.
  \label{eq:expected-increment}
\end{equation}
In particular, if $\beta_t=0$ and
$\E\norm{\bh_t}_{\bH_s}^2\geq\nu_t^2$, then
$\E\cF_{s\to T}\geq\frac12\sum_{t=s+1}^{T}\nu_t^2$.
\end{theorem}

The theorem formalizes gradual forgetting under a falsifiable drift condition.
It does not derive monotonicity from sequential training alone. For $p=1$,
$H_s=1$, $h_{s+1}=1$, and $h_{s+2}=-1$, forgetting rises to $1/2$ and returns
to zero. Under independent innovations, the expected behavior is more explicit.

\begin{corollary}[Independent innovations]
\label{cor:independent-innovations}
Under Theorem~\ref{thm:sequential-accumulation}, suppose
$\{\bh_u\}_{u>s}$ are independent with means $\bm{\mu}_u$ and covariances
$\bSigma_u$. Then
\begin{equation}
  \E\cF_{s\to T}
  =\frac12\norm{\sum_{u=s+1}^{T}\bm{\mu}_u}_{\bH_s}^2
  +\frac12\sum_{u=s+1}^{T}\tr(\bH_s\bSigma_u).
  \label{eq:innovation-accumulation}
\end{equation}
Zero-mean innovation energy accumulates additively in expectation. It grows
linearly in the number of windows only when the per-window trace contribution
has a uniform positive lower bound. In contrast,
coherently aligned mean drift can yield quadratic growth.
\end{corollary}

\begin{proposition}[A quadratic bridge from mixture drift to forgetting]
\label{prop:mixture-bridge}
Let component risks share curvature $\bH\succ0$:
\begin{equation}
  \cL_c(\bx)=\kappa_c+\frac12\norm{\bx-\bm{m}_c}_{\bH}^2,
  \qquad
  \cL_t=\sum_c\pi_{t,c}\cL_c.
  \label{eq:component-quadratic}
\end{equation}
Let $\bm{m}_t=\sum_c\pi_{t,c}\bm{m}_c$, let $\cU$ be a fixed linear PEFT
space, and suppose sequential training reaches the exact constrained optimum
$\bx_t=P_{\cU}^{\bH}\bm{m}_t$, where $P_{\cU}^{\bH}$ is the
$\bH$-orthogonal projector. Then
\begin{equation}
  \cF_{s\to t}
  =\frac12\norm{P_{\cU}^{\bH}(\bm{m}_t-\bm{m}_s)}_{\bH}^2.
  \label{eq:mixture-forgetting}
\end{equation}
Equivalently, with $\bM=[\bm{m}_1,\ldots,\bm{m}_C]$ and
$\Delta\bm{\pi}_u=\bm{\pi}_u-\bm{\pi}_{u-1}$, the sequential increments are
$\bh_u=P_{\cU}^{\bH}\bM\Delta\bm{\pi}_u$. Hence independent zero-mean
mixture innovations accumulate according to
Corollary~\ref{cor:independent-innovations}, while cancelling mixture motion can
restore earlier performance.
\end{proposition}

This proposition derives the drift condition in a stylized shared-Hessian
model; it does not claim that arbitrary mixture drift or nonlinear training
must forget. In particular, drift outside $\cU$ produces acquisition error but
does not move the adapter, separating lack of plasticity from forgetting.

\subsection{Conflict, Fixed-Rank Capacity, and Online Error}

Let $w_s>0$ with $\sum_s w_s=1$ and
$\overline{\cL}_T(\bX)=\sum_s w_s\cL_s(\bX)$. Let $\bX_s^\star$ minimize
the individual risk, $\bX^\star$ minimize the historical mixture risk over the
full update space, $\bX_R^\star$ minimize it over the rank-$R$ class $\cC_R$,
and $\widehat{\bX}_T\in\cC_R$ be the online learner's result.
Define the aggregate per-window-oracle gap
\begin{equation}
  \mathcal{R}_T(\widehat{\bX}_T)
  :=\sum_s w_s[\cL_s(\widehat{\bX}_T)-\cL_s(\bX_s^\star)].
  \label{eq:oracle-retention-gap}
\end{equation}
Unlike Eq.~\eqref{eq:signed-forgetting}, this is a final shared-model gap to
separate per-window global oracles, not checkpoint-based continual-learning
forgetting.

\begin{proposition}[Exact retention-gap decomposition]
\label{prop:error-decomposition}
If the stated global minimizers exist, then
\begin{align}
&\mathcal{R}_T(\widehat{\bX}_T)
  \notag\\
&=\underbrace{\overline{\cL}_T(\bX^\star)
  -\sum_s w_s\cL_s(\bX_s^\star)}_{\mathsf{Conflict}_T}
  +\underbrace{\overline{\cL}_T(\bX_R^\star)
  -\overline{\cL}_T(\bX^\star)}_{\mathsf{Capacity}_{R,T}}
  \notag\\
&\quad+\underbrace{\overline{\cL}_T(\widehat{\bX}_T)
  -\overline{\cL}_T(\bX_R^\star)}_{\mathsf{Online}_{R,T}}.
  \label{eq:three-way-decomposition}
\end{align}
All three terms are nonnegative.
\end{proposition}

The online term includes incomplete history, estimation, allocation, and
optimization error. The capacity term has an exact form in a tractable geometry.

\begin{theorem}[Spectral characterization of rank capacity]
\label{thm:spectral-capacity}
Let $\cC_R=\{\bX:\rank(\bX)\leq R\}$ and suppose
\begin{equation}
  \overline{\cL}_T(\bX)
  =\overline{\cL}_T(\bX^\star)
  +\frac12\norm{
  \bC_{\mathrm{out}}^{1/2}(\bX-\bX^\star)
  \bC_{\mathrm{in}}^{1/2}}_F^2,
  \label{eq:kronecker-quadratic}
\end{equation}
where $\bC_{\mathrm{out}}\succ\bm{0}$ and
$\bC_{\mathrm{in}}\succ\bm{0}$. Defining
$\bY^\star=\bC_{\mathrm{out}}^{1/2}\bX^\star
\bC_{\mathrm{in}}^{1/2}$ gives
\begin{equation}
  \mathsf{Capacity}_{R,T}
  =\frac12\sum_{j>R}\sigma_j^2(\bY^\star),
  \qquad
  \bX_R^\star=\bC_{\mathrm{out}}^{-1/2}
  [\bY^\star]_R\bC_{\mathrm{in}}^{-1/2}.
  \label{eq:spectral-capacity}
\end{equation}
\end{theorem}

\begin{corollary}[Global allocation across separable layers]
\label{cor:global-layer-allocation}
Suppose the local historical loss separates over $L$ adapted matrices as
\begin{equation}
  \overline{\cL}_T(\{\bX_\ell\})
  =\overline{\cL}_T(\{\bX_\ell^\star\})
  +\frac12\sum_{\ell=1}^{L}
  \norm{\bC_{\mathrm{out},\ell}^{1/2}
  (\bX_\ell-\bX_\ell^\star)
  \bC_{\mathrm{in},\ell}^{1/2}}_F^2,
  \label{eq:layer-separable-quadratic}
\end{equation}
with positive-definite factors. Under the global budget
$\sum_\ell\rank(\bX_\ell)\leq R$, an optimum retains the $R$ largest values
in the multiset of all whitened singular values
$\{\sigma_j(\bY_\ell^\star)\}_{\ell,j}$, where
$\bY_\ell^\star=\bC_{\mathrm{out},\ell}^{1/2}
\bX_\ell^\star\bC_{\mathrm{in},\ell}^{1/2}$. Its capacity error is one half
the sum of squares of the unselected values.
\end{corollary}

\begin{corollary}[Finite-budget capacity lower bound]
\label{cor:capacity-lower-bound}
Under Theorem~\ref{thm:spectral-capacity}, suppose $\bY^\star$ has $m$
singular values at least $\alpha>0$. Every rank-$R$ adapter with $R<m$ then
incurs
\begin{equation}
  \mathsf{Capacity}_{R,T}\geq\frac12(m-R)\alpha^2.
  \label{eq:capacity-lower-bound}
\end{equation}
For a time-indexed optimum, define
$m_T(\alpha):=\#\{j:\sigma_j(\bY_T^\star)\geq\alpha\}$. At any time $T$, the
capacity lower bound is
$\frac12\max\{m_T(\alpha)-R,0\}\alpha^2$. This bound grows
linearly in $m_T(\alpha)-R$; the realized capacity error need not be monotone
in time because modes may cancel or lose energy.
\end{corollary}

This is the precise capacity result: if independent curvature-weighted modes
continue to enter $\bY^\star$, a fixed $R$ leaves a growing spectral tail. It
also proves that the curvature-whitened truncated SVD in
Eq.~\eqref{eq:weighted-svd}---the per-window solve \method{} runs---is the
locally optimal rank-$R$ operator under the same assumptions, and hence a locally
optimal merge. It does not establish optimality for general nonlinear model
merging.

Capacity is not the whole story. With
$\cL_1(x)=\frac12(x-1)^2$ and
$\cL_2(x)=\frac12(x+1)^2$, the shared optimum $x=0$ has positive conflict even
though a one-dimensional model has no capacity error. Increasing rank cannot
make one unconditional scalar output equal both $1$ and $-1$.

\subsection{The Locally Safe Space Contracts}

Let $\bK_s\succeq0$ be a Gauss--Newton or Fisher sensitivity surrogate in
fixed ambient adapter coordinates, and define
\begin{equation}
  \bA_0^{\mathrm{amb}}=\bm{0},
  \qquad
  \bA_t^{\mathrm{amb}}=\sum_{s=1}^{t}\omega_s\bK_s,
  \qquad \omega_s>0.
  \label{eq:additive-safe-curvature}
\end{equation}
For a scale-invariant approximate diagnostic and $0\leq\tau<1$, let
$\cS_s^{\mathrm{amb}}(\tau)$ be the span of eigenvectors of $\bK_s$ whose
eigenvalues exceed $\tau\norm{\bK_s}_{\mathrm{op}}$, taking this span to be
$\{0\}$ when $\bK_s=0$.

\begin{theorem}[Contraction of the second-order safe subspace]
\label{thm:safe-subspace}
Let $\cN_0^{\mathrm{amb}}=\R^p$,
$\cN_t^{\mathrm{amb}}=\ker(\bA_t^{\mathrm{amb}})$, and
$\cN_t^{\mathrm{amb},(\tau)}=\cap_{s=1}^{t}
\cS_s^{\mathrm{amb}}(\tau)^\perp$. Then
\begin{equation}
  \cN_t^{\mathrm{amb}}=\bigcap_{s=1}^{t}\ker(\bK_s)
  \subseteq\cN_{t-1}^{\mathrm{amb}},
  \qquad
  \cN_t^{\mathrm{amb},(\tau)}
  \subseteq\cN_{t-1}^{\mathrm{amb},(\tau)}.
  \label{eq:safe-kernel}
\end{equation}
The approximate construction is invariant to a positive rescaling of any
single $\bK_s$. Moreover, within an allocation epoch let the fixed dictionary
$\bB\in\R^{p\times m}$ have full column rank, define
$\bG_s=\bB^\top\bK_s\bB$ and
$\bA_t^{\bB}=\bB^\top\bA_t^{\mathrm{amb}}\bB$. Its coefficient-space safe set
satisfies
\begin{equation}
  \cN_t^{\bB}:=\ker(\bA_t^{\bB})
  =\{\ba:\bB\ba\in\cN_t^{\mathrm{amb}}\}
  =\bigcap_{s=1}^{t}\ker(\bG_s)
  \subseteq\cN_{t-1}^{\bB}.
  \label{eq:coefficient-safe-kernel}
\end{equation}
\end{theorem}

This proves a global contraction statement in fixed ambient local coordinates
and, within a fixed dictionary, in atom coefficients. As independent
historical sensitivity directions accumulate, neither the exact nullspace nor
the conservative thresholded safe space can expand. Repeating the same
sensitive subspace does not reduce either dimension merely by increasing its
curvature scale.

The coefficient result compares times only while $\bB$ is fixed. An invertible basis change
$\bB'=\bB\bT$ requires the congruence transport
$\bA_t^{\bB'}=\bT^\top\bA_t^{\bB}\bT$; a rank-reducing re-solve is not invertible.
\method{} sidesteps this by never carrying coefficient state across windows: each
window re-solves Eq.~\eqref{eq:weighted-svd} afresh from the bounded buffer, so
the relevant contraction is the ambient one of Eq.~\eqref{eq:safe-kernel}, which
remains well defined. The coefficient theorem is used only within a single
solve.

The next result connects safe-space exhaustion to plasticity.

\begin{proposition}[Price of a useful update]
\label{prop:price-useful-update}
Let $\bA\succeq0$ be accumulated historical curvature, let $\bz\neq0$ be a
current linearized gain direction, and let $\gamma>0$. For
\begin{equation}
  \inf_{\ba}\frac12\ba^\top\bA\ba
  \quad\text{s.t.}\quad \bz^\top\ba\geq\gamma,
  \label{eq:price-problem}
\end{equation}
the minimum is zero if
$\Pi_{\ker(\bA)}\bz\neq0$. If
$\bz\in\range(\bA)\setminus\{0\}$, it equals
\begin{equation}
  \frac{\gamma^2}{2\bz^\top\bA^\dagger\bz},
  \quad\text{attained at}\quad
  \ba^\star=\frac{\gamma\bA^\dagger\bz}
  {\bz^\top\bA^\dagger\bz}.
  \label{eq:price-solution}
\end{equation}
\end{proposition}

Thus a new window can obtain locally free gain while its descent direction has
a component in the historical nullspace. Once that component disappears, any
fixed gain incurs positive predicted old loss. If the nullspace projection is
arbitrarily small, however, the zero-cost construction can require a large
step outside the local neighborhood; a practical trust-region radius must then
be imposed.

\subsection{Optimal Partial-Rank Allocation in a Block Geometry}

\begin{theorem}[Optimal top-$k$ screening for a stationary block surrogate]
\label{thm:topk-allocation}
Partition coefficients into $M$ fixed atom blocks, let $1\leq k\leq M$, and
define the damped block curvature
$\widetilde{\bA}_j=\bA_j+\lambda\bI\succ0$. Suppose the screening cost and
linearized current gain are
\begin{equation}
  \widetilde Q(\ba)=\frac12\sum_{j=1}^{M}
  \ba_j^\top\widetilde{\bA}_j\ba_j,
  \qquad
  G(\ba)=\sum_{j=1}^{M}\bz_j^\top\ba_j.
  \label{eq:block-problem}
\end{equation}
Let $q_j=\bz_j^\top\widetilde{\bA}_j^{-1}\bz_j$. This subproblem assumes that any
historical linear term has vanished at projected stationarity or is enforced
separately. Let $\cS_k$ contain up to $k$ blocks with the largest strictly
positive scores and assume $\sum_{j\in\cS_k}q_j>0$. Among updates using at most
$k$ blocks and satisfying $G(\ba)\geq\gamma>0$, $\cS_k$ is optimal, with
\begin{equation}
  \widetilde Q_{\min}=\frac{\gamma^2}{2\sum_{j\in\cS_k}q_j},
  \qquad
  \ba_j^\star=\frac{\gamma\widetilde{\bA}_j^{-1}\bz_j}
  {\sum_{\ell\in\cS_k}q_\ell}.
  \label{eq:topk-solution}
\end{equation}
\end{theorem}

If the true positive-definite screening curvature $\widetilde{\bA}$ is
spectrally close to
$\bD=\operatorname{blkdiag}(\widetilde{\bA}_1,\ldots,
\widetilde{\bA}_M)$, namely
$(1-\delta)\bD\preceq\widetilde{\bA}\preceq
(1+\delta)\bD$ with $\delta<1$, the selected
support is within a factor $(1+\delta)/(1-\delta)$ of the true minimum
historical cost (Corollary~\ref{cor:approx-block-allocation}). Without such a
cross-block condition, the top-$k$ rule need not be optimal. This theorem
characterizes optimal partial-rank \emph{screening} for a stationary linear-gain
surrogate; it is the discrete counterpart of the continuous rank-$R$ solve
$\method{}$ performs in Eq.~\eqref{eq:weighted-svd}, where the retained rank is
chosen by truncated SVD rather than block selection. We do not claim the score is
globally optimal beyond this surrogate or for joint LoRA-factor training.

\subsection{Expressive Capacity on a Curvature-Weighted Manifold}
\label{sec:expressivity}

The preceding results measure capacity by rank and by the dimension of a
contracting nullspace. Two adapters with equal rank or equal safe dimension need
not be expressively equivalent: the rank scalar drops the direction of the
attained subspace. We now make ``how much of the update space is actually
occupied, and relative to which historical demand'' into a single geometric
object. A fixed-rank LoRA atom pool and a full fine-tuning update can then be
placed on the same curvature-weighted axis as an \emph{empirical} comparison
through their effective ranks; we treat this correspondence as a comparator, not
as a theorem that unifies the two training regimes.

Let $\bA_t^{\mathrm{amb}}=\sum_{s\le t}\omega_s\bK_s$ be the accumulated
historical curvature from Eq.~\eqref{eq:additive-safe-curvature} and let
$\cU\subseteq\R^p$ be the subspace spanned by the currently active update
directions of either a rank-$R$ LoRA pool or a full fine-tuning update with an
effective rank $r_{\mathrm{eff}}=\dim(\cU)$. The curvature-whitened embedding
$\phi_t(\cU)=(\bA_t^{\mathrm{amb}}+\lambda\bI)^{1/2}\cU$ sends each occupied
subspace to a point of the Grassmann manifold, and the principal-angle distance
\begin{equation}
  d_t(\cU,\cV)
  =\Big(\sum_i\sin^2\theta_i\big(\phi_t(\cU),\phi_t(\cV)\big)\Big)^{1/2}
  \label{eq:expressivity-distance}
\end{equation}
is a metric on $\mathrm{Gr}(r,p)$ with $d_t(\cU,\cV)=0$ iff the whitened
subspaces coincide (Proposition~\ref{prop:expressivity-metric}). We are careful
about what the whitening does: for $\lambda>0$ the operator is invertible, so
$d_t=0$ holds iff $\cU=\cV$ and the whitening creates no new stream-relative
collapse of distinct subspaces. Its role is quantitative---it reweights the
principal angles between \emph{distinct} subspaces by accumulated historical
curvature, so $d_t$ orders adapters by curvature-weighted overlap. Rank equality
remains necessary but not sufficient for $d_t=0$, which is what lets a
curvature-weighted distance separate equal-rank adapters that a rank scalar
cannot.
The associated effective capacity
\begin{equation}
  \kappa_t(\cU)
  =\frac12\sum_{j>\dim(\cU)}\sigma_j^2
  \big((\bA_t^{\mathrm{amb}}+\lambda\bI)^{1/2}\bX^\star_{\cU^\perp}\big)
  \label{eq:effective-capacity}
\end{equation}
is the curvature-whitened spectral tail of the part of the local optimum
orthogonal to the occupied subspace. Under the Kronecker quadratic model of
Theorem~\ref{thm:spectral-capacity} and \emph{at the optimal rank-$R$ class}
$\cU=\cC_R$, it reduces to the discrete spectral tail
$\frac12\sum_{j>R}\sigma_j^2(\bY^\star)$, recovering
Theorem~\ref{thm:spectral-capacity}. For a generic occupied subspace
$\cU\neq\cC_R$, $\kappa_t(\cU)$ is the residual off $\cU$ and only \emph{lower
bounded} by that optimal tail (Eckart--Young gives a minimum, not an identity;
see Remark~\ref{rem:effective-capacity-optimal-only}); it is not equal to it. We
therefore use the tail identity only at $\cC_R$ and never assert it for an
arbitrary fixed subspace.

This separates two questions the rank scalar conflates: whether two occupied
subspaces are expressively equivalent ($d_t=0$), and how much unoccupied
curvature-weighted capacity remains ($\kappa_t$). As an empirical comparator,
full fine-tuning enters through its effective rank $r_{\mathrm{eff}}<p$ rather
than as $R=\infty$; its expressive capacity is then the tail left by its
effective, not nominal, rank. We do not claim this places full fine-tuning and a
LoRA pool on one axis as a theorem---the rank-$R$ matrix class is a nonlinear
variety, not an $R$-plane, and the two regimes differ in how $\cU$ is
determined---only that the effective-rank comparison is a well-defined empirical
diagnostic. Nor do we claim $d_t=0$ implies output-function equivalence: two
adapters may share a whitened subspace yet place different coefficients on it.
Proofs and the degeneration to Theorem~\ref{thm:spectral-capacity} are in
Appendix~\ref{app:expressivity}.
